% GENERATED FILE. Edit canonical lessons, then run npm run generate:books.
% canonical-source-hash: fnv1a64:58b3d8932c339943
% canonical-lessons: MR-C285-peksha, MR-C285-varshanpasun, MR-R285-comparing-and-time

\chapter{Bigger, Biggest, and How Long}
\label{ch:mr-bigger-biggest-and-how-long}
\hlchaptermodality{285}

\begin{chapteropening}
\textbf{By the end of this chapter:} \emph{I can compare two things with -pekṣā and jāsta, say the most with sarvāt, and say how long with -pāsūn, -paryant and -bhar.}

Complete MR-R285-comparing-and-time: compare things and say how long something has gone on.
\end{chapteropening}

\section[\texorpdfstring{-\mr{पेक्षा} (-pekṣā)}{-पेक्षा (-pekṣā)}]{-\mr{पेक्षा} (-pekṣā) — than, and the most}

\label{lesson:MR-C285-peksha}



\begin{quote}
Say \textquotedblleft{}I want to go to Pune\textquotedblright{}. (\emph{malā puṇyālā jāycaṁ āhe}.) Then \textquotedblleft{}fast\textquotedblright{}. (\emph{jalad}.)
\end{quote}

\subsection*{You'll want to know: -\mr{पेक्षा}, 'than'}
\begin{quote}
\textbf{\mr{गाडी} \mr{बसपेक्षा} \mr{जलद} \mr{आहे}.}
\end{quote}

\begin{quote}
\emph{gāḍī baspekṣā jalad āhe.} — \textquotedblleft{}The train is faster than the bus.\textquotedblright{}
\end{quote}

\textbf{-\mr{पेक्षा}} goes on the thing you compare against, in its bent form: \textbf{\mr{घर}} $\to$ \textbf{\mr{घरापेक्षा}} (\emph{gharāpekṣā}), \textbf{\mr{चहा}} $\to$ \textbf{\mr{चहापेक्षा}} (\emph{cahāpekṣā}). The adjective stays as it is; Marathi needs no \textquotedblleft{}-er\textquotedblright{}. For a quantity, add \textbf{\mr{जास्त}} (\emph{jāsta}, \textquotedblleft{}more\textquotedblright{}) or \textbf{\mr{कमी}}, \textquotedblleft{}less\textquotedblright{}: \textbf{\mr{मला} \mr{चहापेक्षा} \mr{दूध} \mr{जास्त} \mr{आवडतं}} (\emph{malā cahāpekṣā dūdh jāsta āvaḍtaṁ}), \textquotedblleft{}I like milk more than tea\textquotedblright{}.

\subsection*{You'll want to know: \mr{सर्वात}, 'the most'}
\begin{quote}
\textbf{\mr{मीरा} \mr{वर्गात} \mr{सर्वात} \mr{उंच} \mr{आहे}.}
\end{quote}

\begin{quote}
\emph{mīrā vargāt sarvāt uñc āhe.} — \textquotedblleft{}Meera is the tallest in the class.\textquotedblright{}
\end{quote}

\textbf{\mr{सर्वात}} (\emph{sarvāt}) is \textquotedblleft{}of all\textquotedblright{}: \textbf{\mr{सर्वात} \mr{मोठा}} (\emph{sarvāt moṭhā}), \textquotedblleft{}the biggest\textquotedblright{}; \textbf{\mr{सर्वात} \mr{जवळ}} (\emph{sarvāt javaḷ}), \textquotedblleft{}the nearest\textquotedblright{}.

\subsection*{Your turn}
Say these aloud:

\begin{itemize}
\raggedright
  \item \textquotedblleft{}The train is faster than the bus\textquotedblright{} — \emph{gāḍī baspekṣā jalad āhe}
  \item \textquotedblleft{}I like milk more than tea\textquotedblright{} — \emph{malā cahāpekṣā dūdh jāsta āvaḍtaṁ}
  \item \textquotedblleft{}Meera is the tallest in the class\textquotedblright{} — \emph{mīrā vargāt sarvāt uñc āhe}
\end{itemize}

\begin{tcolorbox}[breakable,colback=teal!4,colframe=teal!35!black,title={Before you move on}]
\textquotedblleft{}Than the house\textquotedblright{}? (\textbf{\mr{घरापेक्षा}}.) \textquotedblleft{}More\textquotedblright{}? (\textbf{\mr{जास्त}}.) \textquotedblleft{}The biggest\textquotedblright{}? (\textbf{\mr{सर्वात} \mr{मोठा}}.)
\end{tcolorbox}
\section[\texorpdfstring{\mr{तीन} \mr{वर्षांपासून} (tīn varṣāmpāsūn)}{तीन वर्षांपासून (tīn varṣāmpāsūn)}]{\mr{तीन} \mr{वर्षांपासून} (tīn varṣāmpāsūn) — how long}

\label{lesson:MR-C285-varshanpasun}



\begin{quote}
Say \textquotedblleft{}the train is faster than the bus\textquotedblright{}. (\emph{gāḍī baspekṣā jalad āhe}.) Then \textquotedblleft{}from\textquotedblright{} and \textquotedblleft{}until\textquotedblright{}. (\emph{pāsūn}, \emph{paryant}.)
\end{quote}

\subsection*{You'll want to know: -\mr{पासून}, 'for (and still)'}
\begin{quote}
\textbf{\mr{मी} \mr{तीन} \mr{वर्षांपासून} \mr{पुण्यात} \mr{राहतो}.}
\end{quote}

\begin{quote}
\emph{mī tīn varṣāmpāsūn puṇyāt rāhto.} — \textquotedblleft{}I have lived in Pune for three years.\textquotedblright{}
\end{quote}

Marathi keeps the present: the living still goes on. \textbf{-\mr{पासून}} goes on the stretch of time, in its bent plural: \textbf{\mr{वर्ष}} $\to$ \textbf{\mr{वर्षांपासून}} (\emph{varṣāmpāsūn}), \textbf{\mr{तास}} $\to$ \textbf{\mr{तासांपासून}} (\emph{tāsāmpāsūn}). To ask: \textbf{\mr{तुम्ही} \mr{किती} \mr{दिवसांपासून} \mr{इथे} \mr{आहात}?} (\emph{tumhī kitī divsāmpāsūn ithe āhāt?}), \textquotedblleft{}How long have you been here?\textquotedblright{}

\subsection*{You'll want to know: from … until, and all … long}
\begin{quote}
\textbf{\mr{मी} \mr{सकाळपासून} \mr{संध्याकाळपर्यंत} \mr{काम} \mr{केलं}.}
\end{quote}

\begin{quote}
\emph{mī sakāḷpāsūn sandhyākāḷparyant kām kelaṁ.} — \textquotedblleft{}I worked from morning until evening.\textquotedblright{}
\end{quote}

\textbf{-\mr{पासून} … -\mr{पर्यंत}} frames a stretch with two ends. \textbf{-\mr{भर}}, as in \textbf{\mr{दिवसभर}}, fills it: \textbf{\mr{दिवसभर} \mr{पाऊस} \mr{पडत} \mr{होता}} (\emph{divasbhar pāūs paḍat hotā}), \textquotedblleft{}it rained all day\textquotedblright{}. A bare number of hours needs no ending: \textbf{\mr{मी} \mr{दोन} \mr{तास} \mr{वाचत} \mr{होतो}} (\emph{mī don tās vācat hoto}).

\subsection*{Your turn}
Say these aloud:

\begin{itemize}
\raggedright
  \item \textquotedblleft{}I have lived in Pune for three years\textquotedblright{} — \emph{mī tīn varṣāmpāsūn puṇyāt rāhto}
  \item \textquotedblleft{}How long have you been here?\textquotedblright{} — \emph{tumhī kitī divsāmpāsūn ithe āhāt?}
  \item \textquotedblleft{}I worked from morning until evening\textquotedblright{} — \emph{mī sakāḷpāsūn sandhyākāḷparyant kām kelaṁ}
\end{itemize}

\begin{tcolorbox}[breakable,colback=teal!4,colframe=teal!35!black,title={Before you move on}]
\textquotedblleft{}For three years\textquotedblright{}? (\textbf{\mr{तीन} \mr{वर्षांपासून}}.) Which tense goes with it, for something still going on? (The present: \textbf{\mr{राहतो}}.) \textquotedblleft{}From morning until evening\textquotedblright{}? (\textbf{\mr{सकाळपासून} \mr{संध्याकाळपर्यंत}}.)
\end{tcolorbox}
\section[(review)]{Review — comparing and time}

\label{lesson:MR-R285-comparing-and-time}



\begin{quote}
Say \textquotedblleft{}I got angry\textquotedblright{} and \textquotedblleft{}I want to become a teacher\textquotedblright{}. (\emph{malā rāg ālā}, \emph{malā śikṣak vhāycaṁ āhe}.) Ask what time the Pune train leaves. (\emph{puṇyācī gāḍī kitī vājtā nighte?}) Then \textbf{\mr{ऋतू}} and \textbf{\mr{ज्ञान}}, and the letters \textbf{\mr{ऋ}} and \textbf{\mr{ञ}}.
\end{quote}

\subsection*{Your turn}
Say these aloud:

\begin{itemize}
\raggedright
  \item In my opinion the train is better than the bus.
  \item I like the monsoon season most, with \textbf{\mr{ऋतू}}.
  \item I have lived in Pune for three years.
  \item Yesterday it rained all day. I read from morning until evening.
  \item how long you have been learning Marathi, with \textbf{-\mr{पासून}}
\end{itemize}

\begin{tcolorbox}[breakable,colback=teal!4,colframe=teal!35!black,title={Before you move on}]
\textquotedblleft{}Than\textquotedblright{}? (\textbf{-\mr{पेक्षा}}.) \textquotedblleft{}The most\textquotedblright{}? (\textbf{\mr{सर्वात}}.) \textquotedblleft{}For three years\textquotedblright{}? (\textbf{\mr{तीन} \mr{वर्षांपासून}}.) \textquotedblleft{}All day\textquotedblright{}? (\textbf{\mr{दिवसभर}}.)
\end{tcolorbox}
